\documentclass[11pt,a4paper]{article}

\usepackage[utf8]{inputenc}
\usepackage[T1]{fontenc}
\usepackage{lmodern}
\usepackage[a4paper,top=2.4cm,bottom=2.4cm,left=2.6cm,right=2.6cm]{geometry}
\usepackage{parskip}
\usepackage{enumitem}
\usepackage{xcolor}
\usepackage{amsmath}
\usepackage{booktabs}
\usepackage{fancyhdr}
\usepackage{titlesec}

\definecolor{polito}{RGB}{0,60,113}
\definecolor{rulegray}{RGB}{200,200,200}

\titleformat{\section}{\normalfont\large\bfseries\color{polito}}{\thesection.}{0.5em}{}
\titlespacing*{\section}{0pt}{1.5\baselineskip}{0.45\baselineskip}

\pagestyle{fancy}
\fancyhf{}
\renewcommand{\headrulewidth}{0.4pt}
\fancyhead[L]{\footnotesize Riassunto della tesi --- Dylan Magliano (s337915)}
\fancyfoot[C]{\footnotesize \thepage}

\begin{document}

\begin{center}
  {\Large\bfseries\color{polito} POLITECNICO DI TORINO}\\[0.35em]
  {\normalsize Corso di Laurea Magistrale in Ingegneria Informatica}\\
  {\small\itshape Orientamento Artificial Intelligence and Data Analytics}\\[1.1em]
  {\large\bfseries Dai pixel ai portafogli finanziari: generare scenari\\
   di mercato realistici tramite variational autoencoder}\\[0.7em]
  {\normalsize\bfseries Riassunto della tesi}\\[0.5em]
  {\small Candidato: Dylan Magliano (s337915) \quad---\quad
   Relatori: prof. Flavio Giobergia, dott. Sergio Caprioli,
   dott. Emanuele Cagliero}\\
  {\small Tesi svolta in collaborazione con Intesa Sanpaolo
   \quad---\quad Anno Accademico 2025--2026}
\end{center}

\vspace{0.5em}
\noindent\textcolor{rulegray}{\rule{\textwidth}{0.8pt}}

\section{Contesto e problema di ricerca}

La matrice di correlazione fra i rendimenti degli asset è l'oggetto su cui
poggia gran parte dell'analisi quantitativa dei portafogli, dalla
formulazione di Markowitz in poi. Stimarla dalle sole serie storiche
comporta però due limiti distinti, entrambi strutturali.

Il primo è il \emph{rumore statistico}. Quando il numero di titoli $N$ è dello
stesso ordine di grandezza della lunghezza $T$ della serie usata per la stima,
la matrice empirica diventa instabile e il suo spettro si popola di autovalori
spuri, un fenomeno studiato in modo sistematico dalla Random Matrix Theory.

Il secondo, più profondo, è la \emph{limitatezza campionaria}: il passato
offre una sola traiettoria fra tutte quelle che sarebbero state possibili.
Configurazioni di mercato mai osservate nello specifico campione a
disposizione, ma plausibili sul piano statistico, restano
inaccessibili a qualunque analisi puramente storica.

Da qui la domanda di ricerca della tesi: è possibile costruire un modello
generativo profondo che, appreso il modo in cui le correlazioni si organizzano
nei dati reali, produca matrici di correlazione sintetiche nuove ---
matematicamente valide e strutturalmente indistinguibili da quelle vere ---
con cui popolare uno spazio di scenari più ricco di quello che la storia da
sola può fornire? Il precedente più diretto in letteratura è CorrGAN
(Marti, 2020), che affronta lo stesso problema con reti avversarie; il lavoro
di Caprioli et al.\ (2024) mostra invece l'uso di un Variational Autoencoder
su matrici storiche con ricampionamento nello spazio latente, impostazione
metodologicamente affine a quella qui adottata.

\section{Dati e pre-processing}

Il fondamento empirico è un dataset di prezzi di chiusura rettificati
dell'indice S\&P 500 sul ventennio 2000--2020. Per garantire continuità delle
serie ed evitare distorsioni campionarie, il paniere è stato filtrato
richiedendo quotazioni ininterrotte sull'intero arco temporale: sopravvivono
$N = 362$ titoli. La classificazione settoriale GICS, necessaria alle analisi
di struttura gerarchica, non era disponibile da un provider strutturato ed è
stata ricostruita tramite classificazione automatica affidata a un modello
linguistico di grandi dimensioni (360 titoli su 362 classificati).

Dai prezzi si calcolano i log-rendimenti giornalieri e da questi, con una
procedura a finestre mobili, le matrici di correlazione di Pearson. L'ampiezza
della finestra è fissata a $T_w = 724$ giorni, esattamente il doppio del
numero di asset: la scelta fissa per costruzione il rapporto di spettro
$q = N/T_w = 0.5$, rendendo il livello di rumore omogeneo nel tempo e
analiticamente governabile tramite i limiti di Marchenko--Pastur, e garantisce
che ogni matrice sia di rango pieno e dunque definita positiva. Con un passo
di 10 giorni si ottengono 461 matrici.

Su questo insieme grezzo agisce un filtro di validità strutturale basato sulla
proprietà di Perron--Frobenius empirica, che richiede al primo autovettore
(il \emph{market mode}) di avere componenti tutte dello stesso segno. Le prime
49 finestre non lo soddisfano; l'analisi puntuale mostra che si tratta di
tutte e sole le finestre contenenti osservazioni antecedenti la metà di
ottobre 2001, in cui l'unico titolo aurifero del paniere (NEM) risulta
sistematicamente anticorrelato con il resto del listino nel suo tipico ruolo
di bene rifugio. Non è un artefatto di misura ma la firma di un regime
eccezionale, ed escluderlo delimita esplicitamente il dominio di validità del
modello: restano \textbf{412 matrici valide}, poi mescolate casualmente e
ripartite in 288 / 82 / 42 fra training, validation e test set. Il
mescolamento casuale, in luogo del partizionamento cronologico, neutralizza il
\emph{distribution shift} di regime macroeconomico che nei test preliminari
impediva alle architetture profonde di generalizzare.

\subsection*{Il vincolo di semidefinita positività e la parametrizzazione di Cholesky}

Un passaggio metodologico decisivo riguarda la forma dell'input. Una matrice
di correlazione deve essere semidefinita positiva (PSD): è un vincolo globale
sullo spettro, non imponibile agendo separatamente sui singoli coefficienti.
Nella prima versione della pipeline i modelli ricostruivano direttamente il
triangolo inferiore della matrice di correlazione; le verifiche spettrali
hanno rivelato la comparsa di autovalori negativi, e basta un solo autovalore
negativo --- anche di magnitudo trascurabile --- perché la decomposizione di
Cholesky fallisca e con essa l'intera simulazione di scenari.

La pipeline definitiva risolve il problema alla radice, cambiando l'oggetto
appreso: ogni matrice viene decomposta come $\mathbf{C} = \mathbf{L}\mathbf{L}^T$
e ciò che viene linearizzato e dato in pasto ai modelli è il fattore
triangolare inferiore $\mathbf{L}$, non la matrice. Qualunque vettore il
decoder produca, la matrice $\hat{\mathbf{L}}\hat{\mathbf{L}}^T$ che se ne
ricompone è una matrice di Gram e quindi PSD per costruzione, indipendentemente
dalla qualità dell'apprendimento; una rinormalizzazione finale a diagonale
unitaria riporta il risultato nel dominio della correlazione senza alterare i
segni degli autovalori. L'input canonico è quindi un vettore di
$D = N(N{+}1)/2 = 65\,703$ elementi.

\section{Architetture e protocollo sperimentale}

Il confronto segue un percorso incrementale su quattro modelli, tutti
addestrati sulla medesima parametrizzazione Cholesky e valutati a parità di
dimensione latente $K$:

\begin{enumerate}[leftmargin=1.4em,itemsep=0.25em,topsep=0.35em]
  \item \textbf{Principal Component Analysis}, benchmark lineare che fissa il
        limite massimo di compressione lineare;
  \item \textbf{Autoencoder Lineare}, singolo layer denso senza bias né
        attivazioni, costruito per replicare quel sottospazio e verificare
        empiricamente il teorema di Bourlard e Kamp;
  \item \textbf{Autoencoder Profondo}, design a imbuto
        $2048 \to 1024 \to 512 \to K$ con attivazioni LeakyReLU, che introduce
        la non-linearità;
  \item \textbf{Variational Autoencoder}, che condivide l'encoder a imbuto del
        precedente ma sostituisce il bottleneck deterministico con un modulo
        di inferenza probabilistica ($\boldsymbol{\mu}$,
        $\log\boldsymbol{\sigma}^2$, \emph{reparameterization trick}).
\end{enumerate}

La \emph{loss} di ricostruzione del VAE non è l'usuale errore quadratico
medio, ma la \textbf{Divergenza di Jeffreys} fra le distribuzioni gaussiane
associate alla matrice originale e a quella ricostruita. La motivazione è
precisa: l'MSE tratta i coefficienti come grandezze indipendenti e risulta
dominato dalle correlazioni di ampiezza maggiore, restando di fatto cieco
sugli autovalori più piccoli dello spettro --- proprio quelli che individuano
i portafogli a varianza minima. Le matrici inverse presenti nella formulazione
di Jeffreys ribaltano questa gerarchia e pesano l'errore in termini relativi
lungo ogni direzione dello spettro. Poiché anche il termine di ricostruzione è
così una divergenza, esso è omogeneo al termine di regolarizzazione di
Kullback--Leibler: l'obiettivo impiegato è di tipo VAE, composto dalla
divergenza di Jeffreys nel dominio delle matrici e dalla divergenza di
Kullback--Leibler nello spazio latente, con peso unitario ($\beta = 1$).

La valutazione combina due famiglie di metriche: errori spaziali (MSE, MAE,
norma di Frobenius) e metriche topologiche costruite sul \emph{Minimum
Spanning Tree} (percentuale di archi in comune, indice di Jaccard,
sovrapposizione dei nodi a più alta \emph{betweenness centrality}, scarti di
lunghezza media dei cammini). Le seconde servono a verificare che una
ricostruzione accurata ``in media'' non abbia comunque alterato la gerarchia dei
legami fra i titoli.

\section{Risultati sulla ricostruzione storica}

Le misurazioni confermano empiricamente il teorema di Bourlard e Kamp: PCA e
Autoencoder Lineare convergono, a parità di $K$, sul medesimo limite di
errore. L'Autoencoder Profondo supera nettamente tale soglia a ogni $K$ e
senza mostrare segni di \emph{overfitting}: a $K = 20$ il suo MSE di test
($7.19 \times 10^{-5}$) è di quasi quattro volte inferiore a quello della PCA
($2.73 \times 10^{-4}$). Il vantaggio della non-linearità si rivela di
\emph{efficienza rappresentativa}: l'AE raggiunge già con $K = 10$ la stessa
qualità che la PCA ottiene soltanto con $K = 100$.

Il Variational Autoencoder si colloca in posizione intermedia. A $K = 20$
ottiene un MSE di $1.55 \times 10^{-4}$, \textbf{inferiore} sia a quello della
PCA sia a quello dell'Autoencoder Lineare, pur restando superiore a quello
dell'Autoencoder Profondo deterministico. MAE e norma di Frobenius
restituiscono esattamente la stessa graduatoria, il che indica che
l'ordinamento non è un artefatto della sensibilità dell'MSE agli \emph{outlier}.
Il quadro topologico è ancora più netto e coerente: sull'\emph{edge overlap}
dell'MST il VAE recupera i due terzi della distanza fra modelli lineari e
Autoencoder Profondo (88.4\% contro 83.0\% della PCA e 91.1\% dell'AE). Il
costo che la regolarizzazione variazionale impone alla fedeltà di
ricostruzione è quindi contenuto, e minore di quanto una prima
intuizione suggerirebbe.

Tre run condotte a $K \in \{10, 20, 30\}$, identiche in ogni altro
iperparametro, mostrano un comportamento peculiare del VAE: l'errore di
ricostruzione \emph{peggiora} lievemente al crescere di $K$, all'opposto degli
altri tre modelli. L'analisi dello spettro di attività dello spazio latente ne
suggerisce una spiegazione: in tutte e tre le configurazioni solo circa sei dimensioni
risultano effettivamente attive, mentre le restanti subiscono un
\emph{posterior collapse}. La dimensionalità effettiva della \emph{manifold}
delle correlazioni emerge dunque non come un iperparametro da calibrare, ma come una
proprietà dei dati che il modello individua da sé; aggiungere dimensioni
nominali non accresce la capacità utilizzata e amplia soltanto lo spazio su
cui il vincolo di conformità al \emph{prior} deve essere soddisfatto. La
scelta di $K = 20$ --- compiuta prima delle due run supplementari, sul gomito
della curva di ricostruzione dei modelli deterministici --- non è dunque un
\emph{optimum} empirico del VAE (a $K = 10$ l'errore di ricostruzione è anzi
lievemente inferiore), ma un compromesso metodologico che queste prove
mostrano essere poco critico.

\section{Generazione stocastica di scenari e validazione}

È sul piano generativo che il VAE mostra il proprio valore distintivo.
L'Autoencoder Profondo deterministico, per quanto superiore nella pura
ricostruzione, dispone di uno spazio latente frammentario e privo di
significato al di fuori dei punti osservati: campionarvi un vettore casuale
tende a produrre output privi di coerenza economico-finanziaria. La regolarizzazione
di Kullback--Leibler, unita al rumore di campionamento, modella invece uno
spazio latente continuo e denso, in cui è possibile campionare in modo
significativo. La sintesi più efficace dei risultati è questa complementarità:
\emph{l'AE per ricordare, il VAE per immaginare}.

La procedura di generazione effettivamente adottata non è però un semplice
campionamento dal \emph{prior} $\mathcal{N}(\mathbf{0},\mathbf{I})$, che
produrrebbe uno scenario plausibile qualsiasi, scollegato dal presente.
L'obiettivo è più specifico: costruire una distribuzione di scenari sintetici
della struttura di correlazione \emph{a un orizzonte nominale di 21 giorni a
partire dallo stato di mercato osservato oggi}, con cui stimare la possibile
variabilità delle correlazioni in ottica di rischiosità.
La tecnica impiegata è un \textbf{block-bootstrap storico nello spazio
latente}, articolato in quattro passaggi: l'encoder ricostruisce la traiettoria
latente storica giorno per giorno; se ne calcolano gli incrementi giornalieri;
si campiona casualmente un blocco di 21 incrementi consecutivi (un mese
lavorativo) e se ne cumula il salto; il salto viene sommato allo stato attuale
e decodificato, ottenendo per costruzione una matrice PSD. Ripetendo mille
volte si ottiene una distribuzione di evoluzioni possibili, ancorata al
presente ma diversificata dalla variabilità storica campionata a blocchi.

I mille scenari così prodotti sono stati validati rispetto ai cinque fatti
stilizzati delle matrici di correlazione finanziarie, ciascuno tradotto in un
indicatore scalare confrontabile con la matrice odierna e con l'intervallo
osservato sulle 412 matrici reali:

\begin{enumerate}[leftmargin=1.4em,itemsep=0.2em,topsep=0.35em]
  \item \textbf{Distribuzione \emph{pairwise} asimmetrica verso valori
        positivi}: media fra 0.4609 e 0.5352 su tutti gli scenari, contro
        0.4512 della matrice odierna.
  \item \textbf{Spettro e Marchenko--Pastur}: autovalore di mercato
        $\lambda_1$ nell'intervallo $[172.4,\ 197.2]$ contro 169.1 odierno,
        circa due ordini di grandezza oltre il limite superiore del
        \emph{bulk}, accompagnato da 7--9 modi settoriali.
  \item \textbf{Proprietà di Perron--Frobenius}: componente minima del
        \emph{market mode} sempre positiva, con margine dalla soglia di
        violazione superiore a quello della stessa matrice odierna.
  \item \textbf{Struttura gerarchica settoriale}: rapporto fra correlazione
        media intra-settore e inter-settore sempre maggiore di 1 e interno
        all'intervallo storico; i dendrogrammi riproducono visivamente gli
        stessi blocchi cromatici settoriali della matrice reale.
  \item \textbf{Topologia \emph{scale-free} del MST}: esponente della legge di
        potenza compreso fra 2.053 e 2.306, dentro l'intervallo
        $2 < \gamma < 3$ atteso dalla letteratura.
\end{enumerate}

L'esito è che \textbf{tutti e 1000 gli scenari soddisfano simultaneamente i
cinque criteri}, e tutte le matrici risultano definite positive, in linea con la
semidefinita positività garantita per costruzione dalla parametrizzazione
di Cholesky. L'ispezione grafica di
dettaglio su un singolo scenario corrobora il quadro anche sugli oggetti che
nessun indicatore scalare può riassumere: la similarità coseno rispetto alla
matrice odierna è 0.9999 sul \emph{market mode} e 0.9893 e 0.9883 sui due modi
settoriali successivi.

\section{Limiti e sviluppi futuri}

Il lavoro dichiara esplicitamente i propri limiti metodologici. La verifica
aggregata riduce ciascun fatto stilizzato a un indicatore scalare, mentre
l'ispezione visiva completa resta limitata a un singolo scenario
rappresentativo; i mille scenari derivano inoltre da blocchi di 21 giorni
ampiamente sovrapposti, con numerosità campionaria effettiva inferiore a
mille, e l'intera verifica è ancorata a un unico stato di mercato di partenza.
Il mescolamento casuale, pur affrontando il \emph{distribution shift}, non
elimina la sovrapposizione informativa fra finestre adiacenti (fino al 98.6\%
di osservazioni condivise), il che costituisce un rischio limitato ma reale di
ottimismo nelle metriche di generalizzazione. Il confronto fra modelli, infine,
non copre uniformemente tutte le dimensioni latenti per il VAE, e la
classificazione settoriale GICS non è stata validata con fonti ufficiali.

Il proseguimento naturale è l'applicazione della distribuzione di scenari
generata alla gestione del rischio di portafoglio, confrontando il quadro che
se ne ricava con quello ottenibile dai soli dati storici. Seguono il
rafforzamento della validazione statistica --- ripetendo la verifica a partire
da una pluralità di stati iniziali distribuiti lungo il ventennio ed
estendendola da indicatori scalari a test formali di aderenza distributiva ---
il confronto diretto con paradigmi generativi alternativi (CorrGAN, e più in
generale i modelli di diffusione), e l'estensione della pipeline ad altre
\emph{asset class} e a panieri di titoli internazionali.

\section{Conclusione}

La tesi ha progettato, implementato e valutato empiricamente un framework completo per la
compressione, ricostruzione e generazione stocastica di matrici di
correlazione finanziarie realistiche. Il risultato più significativo è che la
parametrizzazione di Cholesky risolve strutturalmente il vincolo di
semidefinita positività, e che su questa base il Variational Autoencoder paga
un costo contenuto in fedeltà di ricostruzione --- restando anzi superiore ai
benchmark lineari --- in cambio della proprietà che rende possibile la
generazione stocastica. La conformità simultanea dei mille scenari generati ai
cinque fatti stilizzati supporta l'impiego dell'architettura proposta come strumento autonomo
per la produzione di ecosistemi finanziari sintetici matematicamente coerenti.

\end{document}
